\chapter{Conclusion}
\label{ch:sum}

This thesis has presented the design, implementation, and empirical evaluation
of a locally deployed retrieval-augmented generation chatbot for the ELTE
Faculty of Informatics, enabling students and staff to ask natural-language
questions about curriculum, prerequisites, exam rules, and Neptun procedures
without sending any data to an external server.

% ─────────────────────────────────────────────────────────────
\section{Summary of Contributions}
\label{sec:sum-contributions}

\paragraph{Document collection and ingestion pipeline.}
A resumable breadth-first crawler and an incremental ingestion pipeline with
SHA-256 change detection built a 6112-chunk ChromaDB index from 482 faculty
documents, ensuring that subsequent runs only reprocess new or modified files.

\paragraph{Offline RAG pipeline.}
Query responses are grounded in the index using \texttt{all-MiniLM-L6-v2}
embeddings and top-5 cosine retrieval. The entire pipeline runs on a standard
laptop without GPU acceleration or any external network dependency at inference
time.

\paragraph{Local language model serving.}
Two quantised local models --- \texttt{llama3.2:3b} and \texttt{gemma3:4b} ---
are served via Ollama at 4-bit quantisation. No query, retrieved chunk, or
generated answer leaves the host machine during inference.

\paragraph{Production-quality backend and frontend.}
A FastAPI backend exposes the pipeline through a REST API, logs interactions to
a local SQLite database, and serves a browser-based chat interface matching the
faculty's visual identity. A 16-case end-to-end test suite verified the full
request path from API layer to Ollama inference; all tests passed.

\paragraph{Empirical evaluation.}
A 20-question benchmark across four configurations confirmed that RAG is the
decisive factor in answer quality. ROUGE-L improves by $2.5\times$ for
\texttt{llama3.2:3b} and $4.7\times$ for \texttt{gemma3:4b} when retrieval
context is added. The best configuration, \texttt{gemma3:4b} with RAG, achieves
ROUGE-L 0.357, semantic similarity 0.778, and a perfect out-of-scope refusal
rate of 100\%.

% ─────────────────────────────────────────────────────────────
\section{Reflection on Goals}
\label{sec:sum-reflection}

The system meets its two primary goals. All computation runs on the local
machine with no external network dependency after the initial model download,
satisfying the offline-operation requirement. With RAG, both models produce
answers substantially more faithful to faculty source documents than their
no-RAG baselines, which have no parametric knowledge of ELTE-specific rules.
The retrieval step achieves a hit-rate of 0.933, confirming that the vector
index reliably surfaces relevant content.

The 30-second latency target (NF-4) was not met on CPU-only hardware. Mean
response times range from 41.5 to 66.9 seconds across the two RAG
configurations, and \texttt{gemma3:4b} without RAG reaches 157 seconds ---
effectively unusable. This is a hardware constraint rather than an
architectural failure, and is addressed in the future work below.

% ─────────────────────────────────────────────────────────────
\section{Limitations}
\label{sec:sum-limitations}

\paragraph{Retrieval ceiling.}
The hit-rate of 0.933 leaves one in fifteen questions without the most relevant
chunk. The single miss --- Q14, on Stipendium Hungaricum and Erasmus
eligibility --- suggests the chunking strategy does not represent shallow
navigation text well. A cross-encoder re-ranking step is the most direct
remedy.

\paragraph{Out-of-scope refusal reliability.}
The prompt-level refusal mechanism has a known gap: \texttt{llama3.2:3b} with
RAG answered Q19 because a population figure appeared incidentally in a
retrieved chunk, providing apparent grounding for an off-topic response. A
retrieval-score threshold below which the system returns a fixed refusal would
close this gap more robustly than prompt engineering alone.

\paragraph{Benchmark scale.}
The 20-question test set is skewed toward prerequisite rules and credit
requirements. Topics such as exam appeals, credit transfer, and dormitory
applications are not represented, limiting the generalisability of the
evaluation results.

% ─────────────────────────────────────────────────────────────
\section{Future Work}
\label{sec:sum-future}

\paragraph{University-wide deployment.}
The system is architecturally ready for deployment beyond a single laptop. On
shared faculty hardware with a mid-range GPU, a larger 7B-parameter model
would bring latency well within the 30-second target while improving answer
quality. Deployment at the university level would also raise the question of
corpus governance: the current SQLite-backed document store and ChromaDB index
would need to be replaced or augmented with a managed database that supports
access control, scheduled re-crawls, and version-tracked document updates. A
lightweight administration interface for adding new sources, retiring outdated
pages, and triggering incremental re-indexing would be a prerequisite for
sustainable operation at that scale.

\paragraph{Retrieval-score threshold for refusal.}
Comparing the cosine similarity of the top retrieved chunk against a calibrated
threshold would provide a robust, model-independent refusal mechanism,
eliminating the inconsistency observed with \texttt{llama3.2:3b} and the
incidental-match edge case that caused it to answer Q19.

\paragraph{Improved retrieval.}
A cross-encoder re-ranker or hybrid BM25 + dense retrieval would improve
precision on exact-match queries such as course codes and regulation numbers
that embedding models handle less reliably than semantic queries.

\paragraph{Multilingual support.}
Replacing \texttt{all-MiniLM-L6-v2} with a multilingual sentence encoder would
extend the system to Hungarian-language queries without re-indexing the corpus.

\paragraph{User feedback loop.}
The SQLite interaction log provides a natural foundation for a feedback
mechanism. Aggregated helpfulness ratings could identify retrieval failures,
surface low-quality chunks, and prioritise document updates.